%%%%%%%%%%%%%%%%%
% CV tailored for SOFTEAM - Responsable IVV
% Positioning: Ingénieur IVV Système / Responsable IVQ | intégration fonctionnelle, simulation, anomalies & coordination multi-métiers
%%%%%%%%%%%%%%%%%

\documentclass[8pt,a4paper,ragged2e,withhyper]{altacv}

\geometry{left=1.25cm,right=1.25cm,top=.5cm,bottom=1.cm,columnsep=1.2cm}

\usepackage{paracol}
\usepackage{fontawesome5}
\usepackage{hyperref}

\iftutex 
  \setmainfont{Roboto Slab}
  \setsansfont{Lato}
  \renewcommand{\familydefault}{\sfdefault}
\else
  \usepackage[rm]{roboto}
  \usepackage[defaultsans]{lato}
  \renewcommand{\familydefault}{\sfdefault}
\fi

\definecolor{SlateGrey}{HTML}{2E2E2E}
\definecolor{LightGrey}{HTML}{666666}
\definecolor{DarkPastelRed}{HTML}{8AC0D9}
\definecolor{PastelRed}{HTML}{00B0DA}
\definecolor{GoldenEarth}{HTML}{B4AD0C}

\colorlet{name}{black}
\colorlet{tagline}{PastelRed}
\colorlet{heading}{DarkPastelRed}
\colorlet{headingrule}{GoldenEarth}
\colorlet{subheading}{PastelRed}
\colorlet{accent}{PastelRed}
\colorlet{emphasis}{SlateGrey}
\colorlet{body}{LightGrey}

\definecolor{violet}{HTML}{5A2582}
\definecolor{rouge}{HTML}{F53B3B}
\definecolor{noir}{HTML}{00232C}
\definecolor{bleu}{HTML}{00B0DA}
\definecolor{rose}{HTML}{F831A0}
\definecolor{jaune}{HTML}{FFEC00}
\definecolor{vert}{HTML}{34AD6C}

\renewcommand{\namefont}{\Huge\rmfamily\bfseries}
\renewcommand{\personalinfofont}{\footnotesize}
\renewcommand{\cvsectionfont}{\LARGE\rmfamily\bfseries}
\renewcommand{\cvsubsectionfont}{\large\bfseries}
\renewcommand{\cvItemMarker}{{\small\textbullet}}
\renewcommand{\cvRatingMarker}{\faCircle}

\input{../../_shared/pubs-num.tex}

% auto_tex_manager layout tuning
\emergencystretch=3em
\hfuzz=60pt
\hbadness=9999
\sloppy

\begin{document}

\name{BOILEAU Guillaume}
\tagline{Ingénieur IVV Système / Responsable IVQ | Intégration fonctionnelle, simulation, anomalies \& coordination}

\photoL{2.8cm}{../../_shared/moi}

\personalinfo{%
  \email{guillaume.boileaupro@gmail.com}
  \phone{+33 6 59 94 85 84}
  \location{Alpes-Maritimes, FRANCE}
  \linkedin{guillaume-boileau-891b81265}
  \researchgate{Guillaume-Boileau-2}
  \orcid{0000-0002-3576-6968}
}

\makecvheader
\vspace{-0.3cm}

\AtBeginEnvironment{itemize}{\small}
\columnratio{0.64}

\begin{paracol}{2}

\cvsection{Experience}

\cvevent{\textbf{Software Engineer -- Développement logiciel, intégration \& support opérationnel}}{VEV Platform Services France}{2025 -- 2026}{Valbonne, France}

\textbf{Contexte :} Plateforme logicielle industrielle CPMS connectée à des infrastructures de recharge, avec services backend, interfaces applicatives, données opérationnelles, flux d’échanges et contraintes d’exploitation.

\textbf{Objectif :} Développer, intégrer et fiabiliser des services logiciels opérationnels, en intervenant sur les composants applicatifs, les interfaces de données, les incidents techniques, la validation des comportements et la documentation.

\begin{itemize}
  \item \textbf{Intégration logicielle :} Analyse des interfaces entre services applicatifs, données opérationnelles, équipements connectés et contraintes terrain.
  \item \textbf{Support opérationnel :} Diagnostic d’incidents, analyse des causes, formalisation de constats et contribution aux actions de correction.
  \item \textbf{Validation opérationnelle :} Vérification de la cohérence entre données applicatives, comportement des équipements et attendus fonctionnels.
  \item \textbf{Analyse de flux :} Investigation d’échanges FTP, interruptions de flux, incohérences de données et comportements inattendus.
  \item \textbf{Développement applicatif :} Contribution à des composants TypeScript, Angular et Node.js intégrés dans une plateforme opérationnelle.
  \item \textbf{Documentation technique :} Rédaction de constats, analyses d’anomalies et éléments utiles à la résolution par les équipes techniques.
  \item \textbf{Outils \& collaboration :} Utilisation de Git, Linux, Zenhub, documentation technique et suivi collaboratif des sujets.
\end{itemize}

\divider

\cvevent{\textbf{Ingénieur R\&D senior -- IVV, segment sol LISA \& validation système}}{CNES}{2023 -- 2025}{Nice, France}

\textbf{Contexte :} Développement logiciel et activités d’IVV pour le segment sol de la mission spatiale LISA ESA/NASA, avec chaînes de données mission L0--L3, simulation, intégration de modèles, validation technique et environnement CNES/ESA.

\textbf{Objectif :} Mettre en place des moyens logiciels et numériques permettant d’intégrer, vérifier, valider et comparer des contributions hétérogènes dans le cadre du cycle en V et des chaînes de données mission.

\begin{itemize}
  \item \textbf{IVV / IVQ système :} Définition de contrôles de cohérence, règles de comparaison, critères d’acceptation et procédures de vérification.
  \item \textbf{Moyens de simulation numérique :} Développement de workflows Python permettant de tester, comparer et valider des comportements de modèles dans différents cas d’usage.
  \item \textbf{Segment sol LISA :} Développement de composants logiciels liés aux chaînes de données mission, à l’exploitation scientifique et à la validation des niveaux L0--L3.
  \item \textbf{Cycle en V :} Connaissance approfondie de la logique de spécification, intégration, vérification, validation et qualification appliquée aux produits de données mission.
  \item \textbf{Intégration fonctionnelle :} Consolidation de contributions hétérogènes dans un cadre logiciel commun, traçable et reproductible.
  \item \textbf{Analyse d’anomalies :} Identification d’écarts entre sorties, diagnostic des causes possibles, analyse d’impact et formalisation de faits techniques.
  \item \textbf{Suivi des corrections :} Formalisation de points bloquants, recommandations techniques, suivi de versions et vérification de cohérence après correction.
  \item \textbf{Coordination multi-contributeurs :} Interface avec contributeurs scientifiques, logiciels et institutionnels dans un environnement spatial international.
  \item \textbf{Plateforme Python :} Développement de CATpyLISA pour l’intégration, la comparaison, la validation et la revue technique de modèles.
  \textcolor{DarkPastelRed}{\href{https://gitlab.oca.eu/gboileau/lisa_l3_tools}{\bf GitLab: CATpyLISA}}
\end{itemize}

\divider

\cvevent{\textbf{Data Scientist -- Données capteurs, bruit \& validation de performance}}{Universiteit Antwerpen, Belgium}{2021 -- 2023}{Antwerp, Belgium}

\textbf{Contexte :} Études de performance pour instruments de précision et détecteurs de nouvelle génération, combinant mesures environnementales, bruit, modélisation statistique, configurations instrumentales et validation de résultats.

\textbf{Objectif :} Développer des workflows robustes et reproductibles pour identifier, caractériser et réduire des sources affectant la qualité de mesure et la sensibilité de systèmes complexes.

\begin{itemize}
  \item \textbf{Données capteurs :} Analyse de mesures issues de capteurs environnementaux, corrélations spatiales et effets de propagation.
  \item \textbf{Validation de performance :} Analyse de sensibilité, vérification de cohérence, figures de mérite et évaluation des limites de performance.
  \item \textbf{Traitement du signal :} Identification et caractérisation de contributions de bruit dans des mesures complexes.
  \item \textbf{Configurations système :} Études de différentes configurations instrumentales et de différents niveaux de bruit pour la détection de signaux très ténus.
  \item \textbf{Machine learning :} Approches de réduction de bruits non stationnaires avec canaux témoins, machine learning et deep learning.
  \item \textbf{Reporting technique :} Production de résultats traçables, de synthèses techniques et de supports pour parties prenantes internationales.
\end{itemize}

\divider

\cvevent{\textbf{Ingénieur R\&D junior -- Signal, simulation \& diagnostics techniques}}{ARTEMIS, Observatoire de la Côte d’Azur}{2018 -- 2021}{Nice, France}

\textbf{Contexte :} Activités de R\&D liées à l’analyse de signaux faibles, à la simulation numérique, aux diagnostics de bruit et à la préparation de missions spatiales et d’instruments de détection.

\textbf{Objectif :} Développer des méthodes d’analyse, automatiser des simulations et produire des résultats exploitables dans un environnement scientifique et technique exigeant.

\begin{itemize}
  \item \textbf{Simulation numérique :} Développement d’approches de simulation pour environnements de signaux complexes et jeux de données de grande taille.
  \item \textbf{Analyse de données :} Mise en place de méthodes pour analyser et séparer des signaux faibles et superposés.
  \item \textbf{Validation technique :} Vérification de cohérence, comparaison de résultats, études de sensibilité et limites méthodologiques.
  \item \textbf{Diagnostics :} Analyse de sources de bruit, corrélations entre capteurs et risques d’interprétation des résultats.
  \item \textbf{Documentation :} Formalisation des méthodes, hypothèses, résultats et conclusions pour revue collaborative.
\end{itemize}

\switchcolumn

\cvsection{Profile}

\textbf{Ingénieur docteur orienté IVV/IVQ, validation système et intégration fonctionnelle, avec une expérience directe en segment sol LISA, cycle en V, chaînes de données L0--L3, simulation numérique, analyse d’anomalies et coordination technique.}

Positionnement pour ce rôle : mettre en place et exploiter des moyens IVQ, contribuer à l’intégration fonctionnelle, planifier et analyser des essais, investiguer les anomalies, suivre les corrections et coordonner plusieurs métiers dans un contexte système complexe.

\begin{itemize}
  \item Expérience en IVV, validation système, moyens de simulation numérique, contrôles de cohérence et analyse d’écarts.
  \item Connaissance du cycle en V, des chaînes de données mission L0--L3 et de la logique d’intégration / vérification / validation.
  \item Capacité à rédiger des faits techniques, documenter des anomalies, suivre les corrections et produire des synthèses exploitables.
  \item Coordination de contributeurs logiciels, scientifiques et système dans des environnements complexes et internationaux.
  \item Profil en montée sur défense/naval, plateformes à terre, bancs d’essais physiques, sûreté électrique et contraintes RTE.
\end{itemize}

\cvsection{Languages}

\cvskill{English}{4}
\divider
\cvskill{Spanish}{2}
\divider

\medskip

\cvsection{Education}

\cvevent{PhD in Physics}{Université Côte d’Azur}{2018 -- 2021}{Nice, France}

\divider

\cvevent{Selective Graduate Program in Fundamental Physics (Magistère)}{Université Paris-Saclay}{2016 -- 2018}{Orsay, France}

\divider

\cvevent{MSc in Astronomy and Astrophysics}{Observatoire de Paris / Université Paris-Saclay}{2017 -- 2018}{Paris, France}

\divider

\cvevent{BSc in Fundamental Physics}{Université Paris-Saclay}{2015 -- 2016}{Orsay, France}

\cvsection{Professional Training}

\cvevent{Machine Learning in Python with scikit-learn}{FUN MOOC -- Inria}{2026}{France Université Numérique}

\small{
Predictive modelling, supervised learning, model evaluation, cross-validation, metrics, pipelines and practical machine learning workflows with Python and scikit-learn.
}

\divider

\cvevent{Recherche reproductible : principes méthodologiques pour une science transparente}{FUN MOOC -- Inria}{2025}{France Université Numérique}

\small{
Reproducible research, GitLab, notebooks/Jupyter, data management, computational documentation, software environments and reproducibility methods.
}

\divider

\cvevent{Continuous Professional Training -- Software, Data, AI and Systems}{OpenClassrooms}{2024 -- 2026}{}

\small{
Python, SQL, Machine Learning, AI, Linux, JavaScript, TypeScript, Angular, Node.js, Django, REST APIs, C/C++, web development, algorithms and software engineering fundamentals.
}

\end{paracol}

\cvsection{Projects}

\cvevent{CATpyLISA / LISA Segment Sol -- IVV, données L0--L3 et validation}{Mission spatiale LISA ESA/NASA -- CNES/ESA}{2023 -- 2025}{}

Développement logiciel lié au segment sol de LISA : structuration de données mission, niveaux L0--L3, intégration de modèles, comparaison de résultats, validation technique et revue collaborative de workflows scientifiques.

\textbf{Rôle :} développement Python, structuration de données L0--L3, mise en place de moyens de simulation numérique, automatisation d’analyses, validation de modèles, analyse d’écarts, documentation technique et coordination avec plusieurs contributeurs.

\textbf{Compétences mobilisées :} Python, NumPy, SciPy, pandas, Matplotlib, GitLab, segment sol, cycle en V, données L0--L3, IVVQ, data quality, validation de modèles, simulation, documentation.

\divider

\cvevent{Workflows de simulation, performance et analyse de bruit}{Instrumentation scientifique et systèmes complexes}{2018 -- 2025}{}

Développement de workflows numériques et statistiques pour analyser des signaux faibles, caractériser des contributions de bruit, comparer des résultats et évaluer la robustesse de modèles dans des environnements complexes.

\textbf{Rôle :} analyse de performance, simulation, traitement du signal, caractérisation de bruit, études de sensibilité, diagnostics techniques, validation de résultats, reporting et support à la décision.

\textbf{Compétences mobilisées :} traitement du signal, simulation numérique, analyse spectrale, Python, C, validation, analyse de sensibilité, documentation technique.

\divider

\cvevent{VEV -- Intégration logicielle, incidents \& support opérationnel}{Industrial software / Data flows / Connected infrastructure}{2025 -- 2026}{}

Développement et intégration de services logiciels pour une plateforme CPMS connectée à des infrastructures de recharge, avec analyse de flux FTP, investigation d’incidents, vérification de comportements applicatifs, support opérationnel et documentation technique.

\textbf{Rôle :} développement TypeScript, Angular et Node.js, intégration logicielle, analyse de flux FTP, investigation d’incidents, vérification de comportements applicatifs et formalisation de constats techniques.

\textbf{Compétences mobilisées :} TypeScript, Angular, Node.js, FTP, Linux, Git, analyse d’incidents, intégration logicielle, support opérationnel, documentation technique.

\cvsection{Compétences clés}

\vspace{-.3cm}

\cvsubsection{IVV, IVQ \& validation système}

{\small
\cvtag{IVV}
\cvtag{IVQ}
\cvtag{IVVQ}
\cvtag{Intégration fonctionnelle}
\cvtag{Validation système}
\cvtag{Vérification}
\cvtag{Qualification -- ramp-up}
\cvtag{Programmes d’essais}
\cvtag{Analyse d’enregistrements}
\cvtag{Contrôles de cohérence}
\cvtag{Critères d’acceptation}
\cvtag{Cycle en V}
\cvtag{Conformité système}
}

\divider\smallskip
\vspace{-.3cm}

\cvsubsection{Moyens IVQ, simulation \& plateformes}

{\small
\cvtag{Moyens de simulation numérique}
\cvtag{Outillages d’intégration}
\cvtag{Moyens IVQ}
\cvtag{Bancs de tests -- ramp-up}
\cvtag{Plateformes à terre -- ramp-up}
\cvtag{Mise en service -- ramp-up}
\cvtag{Acceptation de moyens -- ramp-up}
\cvtag{Gestion de configuration}
\cvtag{Maintenance préventive -- awareness}
\cvtag{Maintenance corrective -- awareness}
\cvtag{Disponibilité des moyens}
}

\divider\smallskip
\vspace{-.3cm}

\cvsubsection{Anomalies, faits techniques \& corrections}

{\small
\cvtag{Analyse d’anomalies}
\cvtag{Faits techniques}
\cvtag{Suivi de corrections}
\cvtag{Root cause analysis}
\cvtag{Diagnostic technique}
\cvtag{Analyse d’impact}
\cvtag{Investigation}
\cvtag{Remontée de conformité}
\cvtag{Traçabilité}
\cvtag{Documentation technique}
\cvtag{Reporting}
\cvtag{Support décisionnel}
}

\divider\smallskip
\vspace{-.3cm}

\cvsubsection{Coordination multi-métiers \& programme}

{\small
\cvtag{Coordination multi-métiers}
\cvtag{Software}
\cvtag{Systèmes}
\cvtag{Essais}
\cvtag{Énergies -- awareness}
\cvtag{Mécanique -- awareness}
\cvtag{Interface client}
\cvtag{Planning d’intégration}
\cvtag{Jalons programme}
\cvtag{Priorisation}
\cvtag{Rythme d’intégration}
\cvtag{Comptes rendus d’activités}
\cvtag{RTE -- awareness}
\cvtag{Sûreté électrique -- awareness}
}

\divider\smallskip
\vspace{-.3cm}

\cvsubsection{Systèmes complexes, défense/naval \& spatial}

{\small
\cvtag{Systèmes complexes}
\cvtag{Défense/naval -- ramp-up}
\cvtag{Capteurs -- awareness}
\cvtag{CMS -- awareness}
\cvtag{Réseaux -- awareness}
\cvtag{Plateformes -- awareness}
\cvtag{Habilitation défense -- éligible}
\cvtag{Spatial}
\cvtag{Segment sol LISA}
\cvtag{Données mission L0--L3}
\cvtag{CNES}
\cvtag{ESA}
}

\divider\smallskip
\vspace{-.3cm}

\cvsubsection{Logiciel, data \& outils}

{\small
\cvtag{Python}
\cvtag{C}
\cvtag{TypeScript}
\cvtag{Angular}
\cvtag{Node.js}
\cvtag{NumPy}
\cvtag{SciPy}
\cvtag{pandas}
\cvtag{Matplotlib}
\cvtag{Jupyter}
\cvtag{Linux}
\cvtag{Git}
\cvtag{GitLab}
\cvtag{Docker}
\cvtag{CI/CD}
\cvtag{Confluence}
\cvtag{Jira}
\cvtag{Zenhub}
\cvtag{OpenSearch}
\cvtag{Sentry}
}

\end{document}
